Krátké vysvětlení: níže uvedený checklist slouží jako praktický pracovní nástroj pro garanty předmětů při redesignu kurzu s využitím generativní AI.

\begin{enumerate}
  \item[1.] [\ ] Vymezit pedagogické cíle a měřitelné výstupy: formulujte 3–5 klíčových kompetencí.
  \item[2.] [\ ] Zvolit typy hodnocení a jejich vztah k výstupům: preferujte procesní důkazy a u technických úloh kombinaci automatizace a kontroly postupu.
  \item[3.] [\ ] Definovat politiku používání AI v sylabu: specifikujte, co je povoleno, kdy deklarovat AI a případné sankce.
  \item[4.] [\ ] Navrhnout artefakty a metriky hodnocení: zařaďte pole pro deklaraci AI a požadavky na dokumentaci postupu (logy, drafty, verze).
  \item[5.] [\ ] Požadavek na nahrání postupu u cvičení (kde relevantní):
    \begin{itemize}
      \item Vyžadujte nahrání postupu u matematických/výpočetních úloh; některé nástroje to podporují \cite{kb_ai_101_tipu_pdf_04e4a115_page_8_1}.
      \item U rukopisných řešení využijte převod rukopisu (OneNote) pro sběr a automatizaci kontroly \cite{kb_ai_101_tipu_pdf_04e4a115_page_8_1}.
    \end{itemize}
  \item[6.] [\ ] Zvážit inkluzivní a interaktivní aktivity: použijte demonstrace a edukační světy (např. Minecraft Education) \cite{kb_ai_101_tipu_pdf_04e4a115_page_15_2}.
  \item[7.] [\ ] Posoudit rizika ochrany osobních údajů a plán datového toku: identifikujte exportovaná data, potřebu (pseudo)anonymizace a připravte DPA/DPIA podklady.
  \item[8.] [\ ] Pilot a validační plán: pilotujte s malou skupinou; sbírejte data o použitelnosti, přesnosti nástrojů a dopadu na zátěž vyučujících.
  \item[9.] [\ ] Technické a provozní kroky: uložte verze, repozitář, šablony a doporučená nastavení nástrojů a přístupů.
  \item[10.] [\ ] Komunikace se studenty a školení týmu: připravte krátkou informaci o roli AI v kurzu, očekávaných artefaktech a způsobu deklarace.
\end{enumerate}

Doporučení pro rychlé nasazení: před spuštěním pilotu projděte body 5–9 a uložte šablony do online repozitáře (living document). Při práci s interními daty nejprve spusťte kontrolu PII/redakci podle lokálních pravidel a DPA (viz poznámky v Prompt Library) \cite{kb_ai_101_tipu_pdf_04e4a115_page_15_2,kb_ai_101_tipu_pdf_04e4a115_page_8_1}.